\subsection{INVERSE SYSTEMS OF MAGMAS}

DEFINITION 1. An inverse system of magmas relative to the indexing set I is an inverse system of sets \((\mathbf{E}_{\alpha}, f_{\alpha\beta})\) relative to I, each \(E_{\alpha}\) having a magma structure and each \(f_{\alpha\beta}\) being a magma homomorphism.

Let \((\mathbf{E}_{\alpha}, f_{\alpha \beta})\) be an inverse system of magmas whose laws are written multiplicatively. The inverse limit set \(\mathbf{E} = \varprojlim \mathbf{E}_{\alpha}\) is a subset of the product magma \(\prod_{\alpha \in I} \mathbf{E}_{\alpha}\) consisting of the families \((x_{\alpha})_{\alpha \in I}\) such that \(x_{\alpha} = f_{\alpha \beta}(x_{\beta})\) for \(\alpha \leqslant \beta\) . If \((x_{\alpha})\) and \((y_{\alpha})\) belong to E, then for \(\alpha \leqslant \beta\) , \(x_{\alpha} = f_{\alpha \beta}(x_{\beta})\) and \(y_{\alpha} = f_{\alpha \beta}(y_{\beta})\) , hence \(x_{\alpha}y_{\alpha} = f_{\alpha \beta}(x_{\beta})f_{\alpha \beta}(y_{\beta}) = f_{\alpha \beta}(x_{\beta}y_{\beta})\) ; hence E is a submagma of \(\prod_{\alpha \in I} \mathbf{E}_{\alpha}\) . E will be given the law induced by that on \(\prod_{\alpha \in I} \mathbf{E}_{\alpha}\) ; the magma obtained is called the inverse limit magma of the magmas \(E_{\alpha}\) . It enjoys the following universal property:

\begin{enumerate}
\def\labelenumi{(\alph{enumi})}
\item
  For all \(\alpha \in \mathbf{I}\) , the canonical mapping \(f_{\alpha}\) of E into \(\mathrm{E}_{\alpha}\) is a magma homomorphism of E into \(\mathrm{E}_{\alpha} \cdot f_{\alpha} = f_{\alpha \beta} \circ f_{\beta}\) for \(\alpha \leqslant \beta\) .
\item
  Suppose a magma F is given and homomorphisms \(u_{\alpha} \colon \mathrm{F} \to \mathrm{E}_{\alpha}\) such that \(u_{\alpha} = f_{\alpha \beta} \circ u_{\beta}\) for \(\alpha \leqslant \beta\) . There exists one and only one homomorphism \(u \colon \mathrm{F} \to \mathrm{E}\) such that \(u_{\alpha} = f_{\alpha} \circ u\) for all \(\alpha \in \mathrm{I}\) (namely \(x \mapsto u(x) = (u_{\alpha}(x))_{\alpha \in \mathrm{I}}\) ).
\end{enumerate}

If the magmas \(E_{\alpha}\) are associative (resp. commutative), so is E. Suppose that each magma \(E_{\alpha}\) admits an identity element \(e_{\alpha}\) and that the homomorphisms \(f_{\alpha\beta}\) are unital. Then \(e = (e_{\alpha})_{\alpha \in I}\) belongs to E for \(e_{\alpha} = f_{\alpha\beta}(e_{\beta})\) for \(\alpha \leqslant \beta\) and it is an identity element of the magma E; with the above notation, the homomorphisms \(f_{\alpha}\) are unital and if the \(u_{\alpha}\) are unital then u is unital. Further, an element \(x = (x_{\alpha})_{\alpha \in I}\) of E is invertible if and only if each of the \(x_{\alpha}\) is invertible in the corresponding magma \(E_{\alpha}\) and \(x^{-1} = (x_{\alpha}^{-1})_{\alpha \in I}\) ; this follows from the formula \(f_{\alpha\beta}(x_{\beta}^{-1}) = f_{\alpha\beta}(x_{\beta})^{-1} = x_{\alpha}^{-1}\) for \(\alpha \leqslant \beta\) .

From these remarks it can be deduced that if the magmas \(E_{\alpha}\) are monoids (resp. groups) and the \(f_{\alpha\beta}\) are monoid homomorphisms, then the magma E is a monoid (resp. a group). In this case we speak of an inverse system of monoids (resp. groups). The universal property goes over immediately to this case.

It is left to the reader to define an inverse system of rings \((\mathrm{E}_{\alpha}, f_{\alpha \beta})\) and to verify that \(\mathbf{E} = \varprojlim \mathbf{E}_{\alpha}\) is a subring of the product ring \(\prod_{\alpha \in I} \mathbf{E}_{\alpha}\) , called the inverse limit ring of the rings \(\mathbf{E}_{\alpha}\) ; it can be verified that the universal property extends to this case.

Let \(\mathfrak{E} = (\mathrm{E}_{\alpha}, f_{\alpha\beta})\) and \(\mathfrak{E}' = (\mathrm{E}_{\alpha}', f_{\alpha\beta}')\) be two inverse systems of magmas (resp. monoids, groups, rings) relative to the same indexing set. A homomorphism of \(\mathfrak{E}\) into \(\mathfrak{E}'\) is an inverse system \((u_{\alpha})_{\alpha \in I}\) of mappings \(u_{\alpha}: \mathrm{E}_{\alpha} \to \mathrm{E}_{\alpha}'\) such that each \(u_{\alpha}\) is a homomorphism. Under these conditions, the mapping \(u = \varprojlim u_{\alpha}\) of \(\varprojlim \mathrm{E}_{\alpha}\) into \(\varprojlim \mathrm{E}_{\alpha}'\) is a homomorphism (cf.~Set Theory, III, § 7, no. 2).

